% document formatting
\documentclass[10pt]{article}
\usepackage[utf8]{inputenc}
\usepackage[left=1in,right=1in,top=1in,bottom=1in]{geometry}
\usepackage[T1]{fontenc}
\usepackage{xcolor}

% math symbols, etc.
\usepackage{amsmath, amsfonts, amssymb, amsthm}

% lists
\usepackage{enumerate}

% images
\usepackage{graphicx} % for images
\usepackage{tikz}

% code blocks
% \usepackage{minted, listings} 
\usepackage[colorlinks=true, urlcolor=blue]{hyperref}


% verbatim greek
\usepackage{alphabeta}

\graphicspath{{./assets/images/Week 6}}

\title{02-719 Week 6 \\ \large{Genomics and Epigenetics of the Brain}}
 
\author{Aidan Jan}

\date{\today}

\begin{document}
\maketitle

\section*{Review: Single Cell Sequencing}
\begin{itemize}
	\item Each cell is flowed through a tube with a barcoded bead and put into a droplet in oil.  
	\item Enzymes in the droplet with the cell integrate the barcode DNA into the DNA of the cell.
	\item The DNA can then be sequenced.  (And, ideally, we know which cell the DNA is from).
\end{itemize}
We end up with a digital gene expression matrix:
\begin{itemize}
	\item Each cell gets a row
	\item Each gene of interest gets a column
	\item Cells are numbers representing gene expression in each cell
\end{itemize}

\subsection*{Sequencing Process}
\begin{itemize}
	\item After isolating RNA, preparing the cDNA library, and sequencing the cell, we have thousands to millions of reads.
	\item A read is a section of the genome that is read.
	\item Using software, reads can be mapped to locations on the genome.  Areas with a lot of reads imply the gene is expressed, while areas with few reads imply the gene is not expressed.
\end{itemize}

\subsection*{Building the Matrix}
Recall the barcoded bead used in sequencing contains four parts: TruSeq Read 1, Barcode, UMI, and Poly(dT)VN.  
\begin{itemize}
	\item The TruSeq Read 1 gives the gene columns in the matrix
	\item The barcode give the rows of the matrix.
	\item UMI (Unique Molecular Identifier) is used to guarantee that each molecule is only read once.  (After PCR, one cell might be sequenced multiple times, but we only want one.)
\end{itemize}
Afterwards, we need to remove empty droplets (i.e., beads in a droplet with no cell), and \textit{doublets}, droplets containing more than one cell.

\subsubsection*{Resulting Matrix}
The resulting matrix, called the \textbf{DGE}, is the Cell $\times$ Genes matrix we want.  In summary, the steps were
\begin{enumerate}
	\item Alignment
	\item Demultiplexing (Making the rows and columns from the barcode)
	\item UMI Collapse (Removing duplicate UMIs)
	\item Empty Droplet Removal
	\item Ambient RNA Removal (optional)
	\item Doublet Removal (optional)
\end{enumerate}

\subsection*{Properties of the DGE}
\begin{itemize}
	\item It is sparse: most counts are zero
	\begin{itemize}
	    \item Why?  We are sampling very few RNA's per cell compared to the total
    \end{itemize}
	\item It has uneven sampling
	\begin{itemize}
	    \item We are pooling all our cDNA's together when we sequence them so we can't really control the amount that comes from each droplet
    \end{itemize}
\end{itemize}

\subsection*{Correcting for Uneven Sampling: Normalization}
Why normalize:
\begin{itemize}
	\item Allows the expression between cells and genes to be comparable
	\item Makes downstream tasks (i.e., dimensionality reduction, clustering, marker gene finding) a lot simpler
\end{itemize}
\[CPM = \left(\frac{\text{Raw read count for a gene}}{\text{Total mapped reads in the sample}}\right) \times 1000000\]
(CPM = Counts per million)\\\\
For a given cell $i$ and gene $j$, an entry on the matrix is
\[CPM_{i, j} = \frac{\text{Number of UMIs of gene $j$ in cell $i$}}{\text{Number of UMIs in cell $i$}} \times 1000000\]

\subsection*{Normalization for Gene Length}
Longer genes will generate more cDNA reads per molecule of RNA
\begin{itemize}
	\item When comparing genes to each other within a cell, we want to estimate the number of molecules of RNA instead of number of cDNA reads.
\end{itemize}
To normalize, we use TPM (Transcripts per million)
\[TPM_{i, j} = \frac{\text{Number of UMIs of gene $j$ in cell $i$}}{\text{Length of gene $j$}} \times \frac{1000000}{\sum_{\text{All genes $k$}} (\text{Number of UMIs of gene $k$} / \text{Length of gene $k$})}\]

\subsection*{Size Factor Normalization}
\begin{itemize}
	\item Both CPM and TPM make a fairly flawed assumption: The number of molecules of RNA in each cell is equal
	\item This isn't true, for example neurons will express more genes than oligodendrocytes
	\item Size factor normalization makes another assumption: most genes are not significantly differential expressed.
\end{itemize}
For a cell $i$ and gene $j$
\[S_{i, j} = \frac{\text{Number of UMIs of gene $j$ in cell $i$}}{\text{Number of UMIs in gene $j$} / \text{Number of cells}}\]
Line up cell, gene size factors from smallest to largest.  For example,
\[S_{i, 256}, S_{i, 1000}, S_{i, 22}, \dots, S_{i, 252}\]
Genes with the smallest and largest gene size factors are probably differentially expressed.  Those in the middle are probably not differentially expressed.
\begin{itemize}
	\item For a given gene, we define $S_i$ to be the median of size factor$_{i, j}$ for gene $j$.  We use median because large values can skew the mean.
\end{itemize}

\section*{Dimensionality Reduction}
We do this to visualize the dataset.   Clustering and other machine learning is very difficult in high dimensions. (Curse of dimensionality.)  Additionally, we can possibly find features that are more robust than individual genes.

\subsection*{The Metagene Matrix}
\begin{itemize}
	\item Each row is a cell, each column is a metagene.
	\item There are much fewer metagenes than genes.   (This makes the matrix not sparse).
\end{itemize}
When we want to do dimensionality reduction, we do matrix factorization to achieve it.
\begin{center} 
	\includegraphics*[width=0.9\textwidth]{W6_1.png} 
\end{center}
The loadings matrix is the metagene matrix we want.  The components matrix explains how each gene contributes to each metagene.

\subsection*{Principal Component Analysis (PCA)}
PCA exploits the correlation structure in the DGE.
\begin{center} 
	\includegraphics*[width=0.5\textwidth]{W6_2.png} 
\end{center}
\begin{itemize}
	\item Correlated genes will combine to form a single principal component / metagene
	\item These represent molecular pathways (somewhat)
\end{itemize}
PCA is optimized to \textbf{explain the variance} in the matrix.  Principal components are ordered by the amount of variance explained.
\begin{itemize}
	\item If we rotate the graph so that the first $n$ PCs (usually $n = 2$ for visualization) are the graph axes, then all the points have a coordinate based on the two PCs.
	\item The PCs can compress a high dimensional space into much lower (i.e., 2).
\end{itemize}




\end{document}